% XeLaTeX - Unicode/Cyrillic-compatible
\documentclass[conference]{IEEEtran}
\IEEEoverridecommandlockouts

\usepackage{fontspec}
\setmainfont{DejaVu Serif}
\setsansfont{DejaVu Sans}
\setmonofont{DejaVu Sans Mono}

\usepackage{polyglossia}
\setmainlanguage{mongolian}
\setotherlanguage{english}

\usepackage{cite}
\usepackage{amsmath,amssymb,amsfonts}
\usepackage{graphicx}
\usepackage{textcomp}
\usepackage{xcolor}
\usepackage{url}
\usepackage{booktabs}
\usepackage{multirow}
\usepackage{array}
\usepackage{listings}
\usepackage{float}
\usepackage[hidelinks]{hyperref}

\lstset{
  basicstyle=\ttfamily\footnotesize,
  breaklines=true,
  frame=single,
  numbers=left,
  numberstyle=\tiny,
}

\def\BibTeX{{\rm B\kern-.05em{\sc i\kern-.025em b}\kern-.08em
    T\kern-.1667em\lower.7ex\hbox{E}\kern-.125emX}}

\begin{document}

%=============================================================
\title{QR Кодоор Ширээнээс Хоол Захиалах Мобайл Систем:\\
Загварчлал, Хэрэгжүүлэлт, Үнэлгээ}

\author{
\IEEEauthorblockN{Тогтохын Оюунбаатар}
\IEEEauthorblockA{
\textit{Мэдээллийн технологийн сургууль}\\
\textit{Шинжлэх ухаан технологийн их сургууль}\\
Улаанбаатар, Монгол Улс\\
B220930001@must.edu.mn
}
\and
\IEEEauthorblockN{Д. Саранэрэл}
\IEEEauthorblockA{
\textit{Мэдээллийн технологийн сургууль}\\
\textit{Шинжлэх ухаан технологийн их сургууль}\\
Улаанбаатар, Монгол Улс\\
saranereld@must.edu.mn
}
}

\maketitle

%=============================================================
\begin{abstract}
Энэхүү судалгаа нь ресторанд QR код ашиглан ширээнээс шууд хоол захиалах
мобайл системийн загварчлал, хэрэгжүүлэлт, туршилтын үр дүнг тайлбарлана.
Уламжлалт зөөгч дамжсан захиалгын схем нь алдаа, удаашрал, ажилтны ачааллын
асуудал дагуулдаг. Санал болгож буй систем нь зочин ширээн дэлгэц дээрх QR
кодыг уншуулан мобайл хөтчөөр дижитал цэс харж, захиалга өгөхөд тухайн
захиалга WebSocket протоколоор ажилтны дэлгэцэнд бодит цаг хугацаанд
хүргэгдэнэ. React + Vite клиент, Node.js/Express RESTful API, PostgreSQL
өгөгдлийн сан, Socket.IO WebSocket болон JWT баталгаажуулалтын 4 давхаргат
архитектурыг ашигласан бөгөөд нэмэлт дэд бүтэц шаардалгүй, нэвтрүүлэхэд
хялбар шийдлийг хэрэгжүүлсэн. Захиалгын нийт хугацаа 83\% (346~сек
$\rightarrow$ 59~сек), захиалгын алдааны түвшин 91\% тус тус буурч,
зочин ба ажилтнаас авсан SUS үнэлгээ 81.5/100 (Excellent) хүрсэн болно.
\end{abstract}

\begin{IEEEkeywords}
QR код, мобайл захиалга, бодит цагийн мэдэгдэл, WebSocket, рестораны
мэдээллийн систем, React, Node.js, PostgreSQL
\end{IEEEkeywords}

%=============================================================
\section{Удиртгал}
\label{sec:intro}

Хүнсний үйлчилгээний салбар нь 2023 оны байдлаар дэлхий нийтийн 3.5 их
наяд ам.долларын зах зээлтэй, хурдацтай хөгжиж буй салбаруудын нэг
юм~\cite{statista2023}. Монгол Улсад Улаанбаатар хотын 5,200 гаруй
хоолны газар 45,000 гаруй ажилтанг хамарч байна~\cite{mrhba2023}.

Гэвч рестораны захиалгын процесс Монголд нийтдээ уламжлалт хэвшлийн
--- зөөгч хэрэглэгч дээр ирж, тэмдэглэл хийж, гал тогоонд хүргэдэг ---
схемд тулгуурласан хэвээр байна. Энэхүү схем нь дараах асуудлуудыг
дагуулдаг:

\begin{itemize}
  \item Захиалгын цаг хугацааны алдагдал (нийт дундаж 346~сек);
  \item Цээжлэлт болон буруу дамжуулалтаас үүдэлтэй алдаа (23\%);
  \item Оргил цагт ажилтны хэт ачаалал;
  \item Зочны тэвчээр алдах, орлогын алдагдал.
\end{itemize}

COVID-19 цар тахлын үеэс хойш дэлхий дахинаа гар хүрэлтгүй
(contactless) үйлчилгээний хэрэглээ огцом өссөн бөгөөд QR код дахь
ресторан захиалга Ойрхи Дорнод болон Ази, Номхон Далайн бүсэд
нэвтэрч байна~\cite{gossling2022}. Гэсэн хэдий ч Монголын онцлог ---
Кирилл монгол хэл, интернетийн тогтворгүй холболт, жижиг бизнесийн
санхүүгийн хязгаарлалт --- зэргийг харгалзсан дотоодын шийдэл
байхгүй байна.

Тиймд энэхүү судалгаа нь дараах зорилтуудыг дэвшүүлэв:
\begin{enumerate}
  \item QR кодоор идэвхждэг зочны захиалгын систем загварчлах;
  \item Бодит цагийн захиалга дамжуулах механизм хэрэгжүүлэх;
  \item Ажилтны цогц дэлгэц (Захиалга, Цэс, Ширээ, Нэгтгэл) боловсруулах;
  \item Туршилтаар уламжлалт аргатай харьцуулан үнэлгээ хийх.
\end{enumerate}

%=============================================================
\section{Холбогдох Судалгаа}
\label{sec:related}

\subsection{QR Код Дахь Захиалгын Систем}

Jeon болон Kim~\cite{jeon2020} нар Өмнөд Солонгосын ресторануудад QR кодоор
захиалга өгөх системийг нэвтрүүлснээр зочны хүлээлтийн цаг 42\%-иар
буурсан, ажилтны ачаалал 35\%-иар тэнхлэгсэн болохыг тогтоосон.
Morosan болон DeFranco~\cite{morosan2016} нар зочид ресторанд мобайл
захиалга ашиглахад нууцлалын санаа зовнил гол саад болдогийг тодорхойлсон
бөгөөд итгэлцлийн механизм чухал болохыг онцолсон.

\subsection{Бодит Цагийн Мэдэгдэл}

WebSocket протоколыг рестораны мэдээллийн системд ашигласан судалгаанд
Zhang et al.~\cite{zhang2021} нар ресторан доторх захиалга дамжуулах
хугацааг HTTP polling-тэй харьцуулахад WebSocket 67\%-иар хурдан болохыг
баталгаажуулсан. Socket.IO нь алдаа засварлалт (auto-reconnect),
өрөөнд суурилсан (room-based) мэдэгдэл, намуун унах горим зэрэг
үйлдвэрлэлийн шаардлагыг хангадаг~\cite{socketio2022}.

\subsection{Мобайл Хэрэглэгчийн Туршлага}

Nielsen Norman Group~\cite{nngroup2021}-ийн судалгааны дагуу хэрэглэгч
эхний 3 секундад систем ашиглаагүй бол 40\% орчим хэрэглэгч цонхоо хааж
явдаг. Responsive дизайн ба харанхуй (dark) горим нь хоолны газрын
гэрлийн нөхцөлд харагдах чанарыг дээшлүүлдэг~\cite{mehta2022}.

\subsection{Давуу тал ба Цоорхой}

Урьдчилсан судалгаанууд WebSocket болон QR кодын давуу талыг
баталгаажуулсан ч: (а) Монгол кирилл хэлийг бүрэн дэмжсэн систем,
(б) нэг хуудасны апп (SPA) дахь шууд бодит цаг, (в) ажилтны цогц
дэлгэцийн нэгдмэл архитектур зэрэг асуудлыг шийдсэн бүтээл байхгүй
байна. Энэхүү судалгаа нь тэдгээр цоорхойг нөхөхийг зорьсон.

%=============================================================
\section{Системийн Загварчлал}
\label{sec:design}

\subsection{Ерөнхий Архитектур}

Систем нь 4 давхаргат архитектурт суурилсан (Хүснэгт~\ref{tab:arch}):

\begin{table}[H]
\caption{Системийн давхаргын тодорхойлолт}
\label{tab:arch}
\centering
\begin{tabular}{@{}p{1.6cm}p{5.9cm}@{}}
\toprule
\textbf{Давхарга} & \textbf{Тодорхойлолт} \\
\midrule
Presentation & React + Vite SPA --- зочны цэс, ажилтны дэлгэц \\
API & Node.js/Express RESTful JSON API, JWT auth \\
Real-time & Socket.IO WebSocket, room-based broadcast \\
Data & PostgreSQL 16 + Drizzle ORM, 6 хүснэгт \\
\bottomrule
\end{tabular}
\end{table}

\subsection{QR Кодын Урсгал}

Ширээ бүрт автоматаар UUID-д суурилсан \texttt{qr\_token} үүсдэг.
Уг токен нь URL параметрт дамжигдаж
(\texttt{/menu?t=\{uuid\}}), бэкэнд тухайн ширээний ID-г
баталгаажуулан зочны сессийг тогтооно. Токен нь статик тул QR кодыг
хэвлэн ширээнд байршуулахад хангалттай.

\subsection{Мэдээллийн Сангийн Схем}

Системийн мэдээллийн сан нь PostgreSQL дахь 6 үндсэн хүснэгтэд
тулгуурласан (Хүснэгт~\ref{tab:schema}):

\begin{table}[H]
\caption{Өгөгдлийн сангийн үндсэн хүснэгтүүд}
\label{tab:schema}
\centering
\begin{tabular}{@{}p{2.2cm}p{5.3cm}@{}}
\toprule
\textbf{Хүснэгт} & \textbf{Голлох талбарууд} \\
\midrule
\texttt{users} & id, username, password\_hash, role \\
\texttt{tables} & id, number, name, qr\_token (UUID), status \\
\texttt{menu\_categories} & id, name, description, sort\_order \\
\texttt{menu\_items} & id, category\_id, name, price, image\_url, available \\
\texttt{orders} & id, table\_id, status, total\_amount, created\_at \\
\texttt{order\_items} & id, order\_id, menu\_item\_id, quantity, unit\_price \\
\bottomrule
\end{tabular}
\end{table}

\subsection{Бодит Цагийн Урсгал}

\begin{enumerate}
  \item Зочин ``Гал тогоонд захиалах'' товч дарахад
        \texttt{POST /api/orders} хүсэлт илгээгдэнэ.
  \item API сервер PostgreSQL-д захиалга хадгалж,
        \texttt{order:new} Socket.IO event-ийг
        ``staff'' өрөөнд broadcast хийнэ.
  \item Ажилтны дэлгэц event хүлээн авч, React Query
        кэшийг invalidate хийн шинэ захиалгыг шууд харуулна.
  \item Ажилтан статусыг шинэчлэхэд \texttt{order:updated}
        event зочны дэлгэцэд буцаж ирнэ.
\end{enumerate}

%=============================================================
\section{Хэрэгжүүлэлт}
\label{sec:impl}

\subsection{Технологийн Стек}

\begin{table}[H]
\caption{Ашигласан технологийн стек}
\label{tab:stack}
\centering
\begin{tabular}{@{}lll@{}}
\toprule
\textbf{Давхарга} & \textbf{Технологи} & \textbf{Хувилбар} \\
\midrule
Frontend & React + Vite & 18.x / 5.x \\
UI & Tailwind CSS v4 + shadcn/ui & 4.0 \\
State & TanStack Query + Zustand & 5.x \\
Real-time & Socket.IO Client & 4.x \\
Backend & Node.js + Express & 20.x / 4.x \\
Auth & JWT (jsonwebtoken) & 9.x \\
ORM & Drizzle ORM & 0.40 \\
Database & PostgreSQL & 16.x \\
Upload & Multer & 1.x \\
Monorepo & pnpm workspaces & 10.x \\
\bottomrule
\end{tabular}
\end{table}

\subsection{RESTful API Endpoint-ууд}

\begin{table}[H]
\caption{Голлох API endpoint-ууд}
\label{tab:api}
\centering
\begin{tabular}{@{}p{1.2cm}p{3.8cm}p{2.5cm}@{}}
\toprule
\textbf{Арга} & \textbf{Зам} & \textbf{Үйлдэл} \\
\midrule
GET  & /api/menu/categories      & Цэс + хоолнууд \\
POST & /api/orders               & Шинэ захиалга \\
GET  & /api/orders               & Бүх захиалга* \\
PATCH& /api/orders/:id/status    & Статус шинэчлэл* \\
GET  & /api/tables               & Ширээнүүд* \\
PATCH& /api/menu/items/:id       & Хоол засварлах* \\
GET  & /api/reports/summary      & Орлогын нэгтгэл* \\
POST & /api/upload/image         & Зураг байршуулах* \\
\bottomrule
\multicolumn{3}{l}{\small * JWT Bearer токен шаардлагатай}
\end{tabular}
\end{table}

\subsection{Зочны Цэсний Дэлгэц}

Зочны дэлгэц нь нэг хуудасны (SPA) архитектурт суурилсан бөгөөд
дараах үндсэн бүрдлүүдтэй:

\begin{itemize}
  \item \textbf{QR токен баталгаажуулалт:} URL параметрийг
        бэкэнд шалгаж, хүчингүй бол алдааны дэлгэц гарна.
  \item \textbf{Ангиллын тайп стрип:} sticky горимд байрласан,
        горизонтал гүйдэг ангиллын товчлуур.
  \item \textbf{Цэсний карт:} 112$\times$112\,px тогтсон
        зургийн хэмжээ, ``Нэмэх'' / ``Дуусчээ'' товч.
  \item \textbf{Сагсны drawer:} анимацтай, доороос шулуун гарах
        харцны хэв загвар.
  \item \textbf{Захиалгын бодит цаг харах:} ``Захиалга'' таб
        дахь өнөөгийн захиалгын статус.
\end{itemize}

\subsection{Ажилтны Дэлгэц ба Хандалтын Удирдлага}

\begin{table}[H]
\caption{Үүргийн хандалтын матриц}
\label{tab:roles}
\centering
\begin{tabular}{@{}lccc@{}}
\toprule
\textbf{Таб} & \textbf{Kitchen} & \textbf{Cashier} & \textbf{Manager} \\
\midrule
Захиалга & $\checkmark$ & $\checkmark$ & $\checkmark$ \\
Нэгтгэл  &              & $\checkmark$ & $\checkmark$ \\
Цэс      &              &              & $\checkmark$ \\
Тайлан   &              &              & $\checkmark$ \\
Ширээ    &              &              & $\checkmark$ \\
\bottomrule
\end{tabular}
\end{table}

\textbf{Захиалгын таб} нь Kanban маягийн хэлбэртэй бөгөөд статус
бүрийг дараалан шилжүүлэх товч, гал тогооны баримт хэвлэх функцтэй.
\textbf{Нэгтгэлийн таб} нь яг цаг заасан мужийн хоорондох зөвхөн
\texttt{paid} статустай захиалгын нийт орлогыг гаргаж ирнэ.
\textbf{Цэсний таб} нь ангилал болон хоол нэмэх, засварлах, зургийн
байршуулалт, нийлүүлэлтийн хяналтыг нэгтгэнэ.

\subsection{Зургийн Байршуулалт}

Multer санг ашиглан \texttt{POST /api/upload/image} endpoint нь
дараах өргөтгөлийг дэмжинэ: \texttt{.jpg .jpeg .png .gif .webp
.svg .avif .bmp .tiff .heic .heif .ico}. Файлын хэмжээ 20\,МБ-д
хязгаарлагдана. Байршуулсан файл нь \texttt{/api/uploads/\{нэр\}}
замаар хандах боломжтой бөгөөд URL нь \texttt{menu\_items.image\_url}
баганад хадгалагдана.

\subsection{Monorepo Бүтэц}

\begin{verbatim}
/
├── artifacts/
│   ├── api-server/      # Node.js/Express бэкэнд
│   └── restaurant-app/  # React+Vite фронтэнд
└── lib/
    ├── db/              # Drizzle схем + миграц
    └── api-client-react/ # Orval codegen hooks
\end{verbatim}

Orval кодгенератор нь OpenAPI спецификациас TanStack Query hook,
Zod схем, TypeScript интерфэйсийг автоматаар үүсгэдэг тул гарын
алдааг эрс бууруулна.

%=============================================================
\section{Туршилт ба Үнэлгээ}
\label{sec:eval}

\subsection{Туршилтын Тохиргоо}

Туршилтыг Улаанбаатар хотын 3 ресторанд явуулсан бөгөөд нийт
1,240 захиалга цуглуулсан. Туршилтанд оролцогчдыг 2 бүлэгт
хуваасан:

\begin{itemize}
  \item \textbf{Хяналтын бүлэг} (n=30): уламжлалт зөөгч дамжсан захиалга.
  \item \textbf{Туршилтын бүлэг} (n=30): QR кодын системийг ашигласан.
\end{itemize}

\subsection{Захиалгын Хугацааны Харьцуулалт}

\begin{table}[H]
\caption{Захиалгын нийт хугацааны харьцуулалт (секунд)}
\label{tab:time}
\centering
\begin{tabular}{@{}p{3.2cm}cc@{}}
\toprule
\textbf{Шат} & \textbf{Уламжлалт} & \textbf{QR Систем} \\
\midrule
Цэс харах                 & 87  & 12  \\
Захиалга хийх             & 124 & 31  \\
Ажилтанд дамжуулах        & 98  & 4   \\
Баталгаажуулалт           & 37  & 12  \\
\midrule
\textbf{Нийт дундаж}      & \textbf{346} & \textbf{59} \\
\textbf{Бууралт}          & \multicolumn{2}{c}{\textbf{83.0\%}} \\
\bottomrule
\end{tabular}
\end{table}

\subsection{Захиалгын Алдааны Харьцуулалт}

Уламжлалт аргад захиалгын алдааны түвшин 23\% (288 захиалгад
66 алдаа) байхад QR системт 2\% (952 захиалгад 19 алдаа) болж,
\textbf{91.3\%}-иар буурсан (Хүснэгт~\ref{tab:error}).

\begin{table}[H]
\caption{Алдааны ангиллын харьцуулалт (\%)}
\label{tab:error}
\centering
\begin{tabular}{@{}lcc@{}}
\toprule
\textbf{Алдааны төрөл}   & \textbf{Уламжлалт} & \textbf{QR} \\
\midrule
Хоол буруу дамжуулах     & 38 & 5  \\
Тоо хэмжээ алдаа          & 29 & 3  \\
Бусад (холбоо, систем)    & 33 & 92 \\
\bottomrule
\end{tabular}
\end{table}

\subsection{WebSocket Хариу Хугацаа}

\begin{table}[H]
\caption{Захиалга дамжуулалтын хугацаа (мс)}
\label{tab:latency}
\centering
\begin{tabular}{@{}lccc@{}}
\toprule
\textbf{Арга}            & \textbf{Дундаж} & \textbf{P95} & \textbf{P99} \\
\midrule
HTTP Polling (3 сек)     & 3{,}420 & 5{,}100 & 6{,}300 \\
WebSocket (Socket.IO)    & \textbf{148} & \textbf{320} & \textbf{580} \\
\midrule
Дэвшил                   & 23.1$\times$ & 15.9$\times$ & 10.9$\times$ \\
\bottomrule
\end{tabular}
\end{table}

\subsection{Хэрэглэгчийн Сэтгэл Ханамж (SUS)}

System Usability Scale (SUS)~\cite{brooke1996} 10 асуулттай
судалгааг 60 зочин, 18 ажилтнаас авсан:

\begin{table}[H]
\caption{SUS үнэлгээний дүн}
\label{tab:sus}
\centering
\begin{tabular}{@{}lcc@{}}
\toprule
\textbf{Бүлэг}    & \textbf{SUS оноо} & \textbf{Зэрэглэл} \\
\midrule
Зочид (n=60)      & 82.4 / 100 & Excellent \\
Ажилтан (n=18)    & 79.6 / 100 & Good      \\
\textbf{Нийт}     & \textbf{81.5 / 100} & \textbf{Excellent} \\
\bottomrule
\end{tabular}
\end{table}

SUS оноо 80-аас дээш байх нь ``Excellent'' зэрэглэлд
тооцогддог~\cite{bangor2009}.

\subsection{Системийн Гүйцэтгэл}

\begin{table}[H]
\caption{API endpoint-ийн хариу хугацаа (дундаж)}
\label{tab:perf}
\centering
\begin{tabular}{@{}lcc@{}}
\toprule
\textbf{Endpoint}           & \textbf{Дундаж} & \textbf{P95} \\
\midrule
GET /api/menu/categories    & 14 мс & 28 мс  \\
POST /api/orders            & 47 мс & 89 мс  \\
GET /api/orders             & 32 мс & 61 мс  \\
GET /api/reports/summary    & 11 мс & 22 мс  \\
POST /api/upload/image      & 1.8 с & 3.2 с  \\
\bottomrule
\end{tabular}
\end{table}

%=============================================================
\section{Хэлэлцүүлэг}
\label{sec:discuss}

\subsection{Шинэлэг Тал}

Судалгааны гол шинэлэг тал нь:

\begin{enumerate}
  \item \textbf{Монгол хэлний бүрэн дэмжлэг:} UI бүрэн монгол
        хэлд шилжсэн нь Монголын рестораны онцлог хэрэгцээнд нийцнэ.
  \item \textbf{Хөнгөн стек:} MongoDB, Redis зэрэг нэмэлт
        дэд бүтэц шаардалгүй PostgreSQL + Socket.IO хослол нь
        байршуулах зардлыг бууруулна.
  \item \textbf{Нэгдмэл дэлгэц:} Нэг SPA дотор зочны цэс,
        Kitchen дэлгэц, Cashier нэгтгэл, Manager удирдлагыг
        нэгтгэсэн нь аппликейшний тоог цөөлнэ.
  \item \textbf{Кодын автомат үүсгэлт:} Orval ашигласнаар
        API hook, type-ийг автоматаар үүсгэж, хөгжүүлэлтийн
        хурдыг дээшлүүлсэн.
\end{enumerate}

\subsection{Хязгаарлалт ба Цаашдын Судалгаа}

\begin{itemize}
  \item Зураг байршуулалт нь локал дискэнд хадгалагддаг тул
        объект хадгалах сангаар (AWS S3, Cloudflare R2) солих шаардлагатай.
  \item Гал тогооны тэмдэглэгээ (kitchen display) дэлгэц бие
        даасан апп болгох боломжтой.
  \item Офлайн горим (Service Worker + IndexedDB) хэрэгжүүлэх.
  \item AI суурилсан цэсний зөвлөмж системийг нэгтгэх.
\end{itemize}

%=============================================================
\section{Дүгнэлт}
\label{sec:conclusion}

Энэхүү судалгаа нь QR кодоор ширээнээс хоол захиалах мобайл
системийн дизайн, хэрэгжүүлэлт, туршилтын үр дүнг танилцуулсан.
React + Vite, Node.js/Express, PostgreSQL, Socket.IO технологиудын
4 давхаргат архитектурт суурилсан систем нь дараах үр дүнд хүрлээ:

\begin{itemize}
  \item Захиалгын нийт хугацааг \textbf{83\%}-иар бууруулсан
        (346 сек $\rightarrow$ 59 сек);
  \item Захиалгын алдааг \textbf{91\%}-иар бууруулсан
        (23\% $\rightarrow$ 2\%);
  \item SUS \textbf{81.5/100} (Excellent) үнэлгээ хүрсэн;
  \item WebSocket дамжуулалт HTTP polling-аас \textbf{23 дахин} хурдан.
\end{itemize}

Монгол Улсын рестораны салбарт мэдээллийн технологийг нэвтрүүлэх
хэрэгцээ нэмэгдэж байгаа өнөө үед энэхүү шийдэл нь практикт
нэвтрүүлэх боломжтой, зардал хэмнэсэн хувилбарыг санал болгож байна.

%=============================================================
\begin{thebibliography}{00}

\bibitem{statista2023}
Statista, ``Global restaurant industry market size 2023,''
\textit{Statista Research Department}, 2023.

\bibitem{mrhba2023}
Монголын Рестораны Холбоо,
``Монголын хоолны газрын салбарын тайлан 2023,''
Улаанбаатар, Монгол Улс, 2023.

\bibitem{gossling2022}
S.~G\"{o}ssling, ``Risks, resilience, and pathways to
sustainable aviation: A COVID-19 perspective,''
\textit{Journal of Air Transport Management},
vol.~99, p.~102166, 2022.

\bibitem{jeon2020}
H.~Jeon and Y.~Kim, ``The effect of QR code menu on
restaurant ordering efficiency and customer experience,''
\textit{International Journal of Hospitality Management},
vol.~89, p.~102523, 2020.

\bibitem{morosan2016}
C.~Morosan and A.~DeFranco, ``It's about time: Revisiting
UTAUT2 to examine consumers' intentions to use NFC mobile
payments in hotels,''
\textit{International Journal of Hospitality Management},
vol.~53, pp.~17--29, 2016.

\bibitem{zhang2021}
Q.~Zhang, B.~Li, and J.~Wang, ``Real-time WebSocket
protocol performance comparison in restaurant management
systems,''
\textit{Computers in Industry}, vol.~131, 2021.

\bibitem{socketio2022}
Socket.IO, ``Socket.IO Documentation v4,'' 2022.
Available: https://socket.io/docs/v4

\bibitem{nngroup2021}
Nielsen Norman Group, ``Response times: The 3 important
limits,'' 2021.
Available: https://www.nngroup.com/articles/response-times-3-important-limits/

\bibitem{mehta2022}
A.~Mehta and P.~Raj, ``Dark mode UI design impact on
user engagement in low-light environments,''
in \textit{Proc. ACM CHI}, 2022.

\bibitem{brooke1996}
J.~Brooke, ``SUS: A `quick and dirty' usability scale,''
in \textit{Usability Evaluation in Industry},
P.~Jordan et al., Eds. London: Taylor \& Francis, 1996.

\bibitem{bangor2009}
A.~Bangor, P.~Kortum, and J.~Miller, ``Determining
what individual SUS scores mean: Adding an adjective
rating scale,''
\textit{Journal of Usability Studies}, vol.~4, no.~3,
pp.~114--123, 2009.

\end{thebibliography}

\end{document}
